\documentclass[11pt,a4paper]{article}
\usepackage[a4paper,margin=24mm]{geometry}
\usepackage[utf8]{inputenc}
\usepackage[T1]{fontenc}
\usepackage{lmodern}
\usepackage{amsmath,amssymb,amsthm,mathtools}
\usepackage{booktabs}
\usepackage{tabularx}
\usepackage{enumitem}
\usepackage{microtype}
\usepackage{xcolor}
\definecolor{linkblue}{rgb}{0,0.15,0.5}
\usepackage[colorlinks=true,linkcolor=linkblue,citecolor=linkblue,urlcolor=linkblue]{hyperref}
\usepackage{url}
\allowdisplaybreaks
\emergencystretch=2em

\renewcommand{\thesection}{S\arabic{section}}
\renewcommand{\theequation}{S.\arabic{equation}}
\renewcommand{\thetable}{S\arabic{table}}
\theoremstyle{plain}
\newtheorem{theorem}{Theorem}[section]
\newtheorem{proposition}[theorem]{Proposition}
\newtheorem{lemma}[theorem]{Lemma}
\newtheorem{corollary}[theorem]{Corollary}
\theoremstyle{definition}
\newtheorem{definition}[theorem]{Definition}
\newtheorem{example}[theorem]{Example}
\theoremstyle{remark}
\newtheorem{remark}[theorem]{Remark}

\newcommand{\R}{\mathbb{R}}
\newcommand{\calI}{\mathcal{I}}
\newcommand{\calH}{\mathcal{H}}
\newcommand{\calK}{\mathcal{K}}
\newcommand{\calV}{\mathcal{V}}
\newcommand{\calS}{\mathcal{S}}
\newcommand{\calF}{\mathcal{F}}
\newcommand{\calP}{\mathcal{P}}
\newcommand{\Pol}{\Pi_{\calI}}
\newcommand{\J}{J_{\calI}(T)}
\newcommand{\Gam}{\Gamma_{\calI}}
\newcommand{\rhoHG}{\rho_{\calH,G}}
\newcommand{\ebar}{\bar e}
\newcommand{\tfr}[2]{\tfrac{#1}{#2}}
\newcommand{\dd}{\,\mathrm{d}}
\DeclareMathOperator{\conv}{conv}
\DeclareMathOperator{\rank}{rank}
\DeclareMathOperator{\TV}{TV}
\DeclareMathOperator{\Lift}{Lift}
\DeclareMathOperator{\depth}{depth}

\title{\textbf{Supplementary Material for\\ ``Finite certificates of nonviability under incomplete observation: a continuous-to-finite bridge with exact witnesses''}}
\author{Amin Abaee\\[0.3em]{\small Independent Researcher}\\[0.2em]{\small ORCID: \href{https://orcid.org/0000-0002-0019-1842}{0000-0002-0019-1842}\quad \href{mailto:amin_abaee@ut.ac.ir}{amin\_abaee@ut.ac.ir}}}
\date{October 15, 2026}

\begin{document}
\maketitle

\noindent
Numbered statements and equations of the main text are cited as ``Theorem~3.2'', ``(10)'', etc.; statements and equations of this file carry the prefix S. All exact values are reproduced by the verification script \path{paper03_computational_certification_v23_verification.py} (63 checks); its output is archived as \path{paper03_computational_certification_v23_verification_output.txt}.

\tableofcontents
\section{Proofs for Sections 2--4 of the main text}\label{S:proofs}

Throughout, labels are $a=(\theta,i,t)\in\calS=\Theta\times\{1,\dots,s\}\times[0,T]$; $k_a$, $\beta_a$, $\calF_a$, $\J$, $\phi_U$, $\Gam$, $\rhoHG$ and program (10) are as in the main text.

\subsection{Proposition 2.2 (robust rows)}
The variation-of-constants formula gives, for $t\le T$,
\[
c_i^\top x_\theta(t)=c_i^\top X_\theta(t,0)x_0+\int_0^tc_i^\top X_\theta(t,\sigma)g_\theta(\sigma)\dd\sigma+\int_0^tk_a(\sigma)^\top u(\sigma)\dd\sigma+\int_0^tc_i^\top X_\theta(t,\sigma)E_\theta(\sigma)d(\sigma)\dd\sigma .
\]
The initial-state term is bounded over $x_0\in P_\theta$ by the support value $h_{P_\theta}(X_\theta(t,0)^\top c_i)$, and the disturbance term, maximised over all measurable $d(\cdot)$ with values in $D_\theta$, equals $\int_0^th_{D_\theta}(E_\theta(\sigma)^\top X_\theta(t,\sigma)^\top c_i)\dd\sigma$; for polytopic $D_\theta$ a maximising vertex selection is measurable, so the bound is attained by an admissible disturbance signal. Collecting terms, $c_i^\top x_\theta(t)\le b_i$ holds for every $(x_0,d)$ if and only if $\calF_a(\pi)\le0$. The separation of the two maximisations from the control term is exactly the statement that the control signal in a mode does not depend on $x_0$ or $d$, which is the policy class of the main text. \qed

\subsection{Proposition 2.3 (value and nonviability)}
The admissible signal space is a finite product of sets $L^\infty(I_r;U)$; because $U$ is compact and convex these are weak-$\ast$ compact (Banach--Alaoglu; closedness because $U$ is a finite intersection of half-spaces and each $\{u:f_i^\top u\le g_i\text{ a.e.}\}$ is weak-$\ast$ closed). Each $\calF_a$ is affine and weak-$\ast$ continuous because $k_a\in L^1$. The supremum of the $\calF_a$ over the finitely many labels of a mesh---and, under (H1), over the compact label set, where it is a supremum of an equicontinuous family---is weak-$\ast$ lower semicontinuous and attains its minimum: $\J$ is well defined. Safety of a policy is $\max_a\calF_a(\pi)\le0$, so $B_0\in\calK^T_\calI$ iff $\J\le0$; the constraints being closed, $\J=0$ admits a safe policy, and nonviability is the strict inequality $\J>0$. \qed

\subsection{Lemma 2.4 (information monotonicity)}
Every $\calI$-adapted policy is $\calI'$-adapted when $\calI'$ refines $\calI$ at every time: the controller with finer information may ignore the extra information. The minimum in (5) is therefore taken over a larger set and $J_{\calI'}(T)\le\J$. \qed

\subsection{Lemma 3.1 (semantics and tightness of the atomic margin)}
Fix finitely many labels $a_\ell$ with weights $\lambda_\ell\ge0$, $\sum_\ell\lambda_\ell=1$. For any policy $\pi$,
\[
\sum_\ell\lambda_\ell\calF_{a_\ell}(\pi)=\sum_{(r,C)}\int_{I_r}g_{r,C}(\sigma)^\top u_{r,C}(\sigma)\dd\sigma-\sum_\ell\lambda_\ell\beta_{a_\ell},
\]
where $g_{r,C}=\sum_{\ell:\theta_\ell\in C}\lambda_\ell k_{a_\ell}$ pools the kernels of the labels whose mode lies in the cell $C$ active on $I_r$ and $u_{r,C}$ is the common control used in that cell; this uses only that the policy is $\calI$-adapted. Pointwise $g^\top u\ge\phi_U(g)=\min_{v\in U}g^\top v$, so the weighted average of the violations is at least $\Gam(\lambda,a)$. If $\Gam(\lambda,a)>0$, the maximum over $\ell$ of $\calF_{a_\ell}(\pi)$ is at least the average and therefore positive: some stored label is violated by at least $\Gam(\lambda,a)$. For tightness, for a.e.\ $\sigma$ the minimisers of $v\mapsto g_{r,C}(\sigma)^\top v$ over the polytope $U$ form a face; choosing the lexicographically first minimising vertex in a fixed vertex order gives a measurable selection $u_{r,C}$ with $g_{r,C}^\top u_{r,C}=\phi_U(g_{r,C})$ a.e. The policy assembled from these selections is $\calI$-adapted and attains the weighted average $\Gam(\lambda,a)$. \qed

\subsection{Proposition 4.1 (arbitrarily large minimal obstructions with a scalar input)}
The construction and the feasibility argument are in the main text; here are the certificate, the rank and the minimality in detail. With $T=q$, $I_j=[j-1,j]$, $u(t)\in[0,1]$ and $w_j=\int_{I_j}u$: mode $j\le q$ has $x_j(T)=w_j-1$ and mode $q+1$ has $x_{q+1}(T)=q-\varepsilon-\sum_jw_j$. The binding rows are the terminal lower bounds $-x\le0$, whose row functionals are
\[
\calF_j(\pi)=1-w_j=\int_0^T\bigl(-\mathbf 1_{I_j}\bigr)u-(-1),\qquad
\calF_{q+1}(\pi)=\sum_jw_j-(q-\varepsilon)=\int_0^T(+1)\,u-(q-\varepsilon),
\]
so the kernels are $k_j=-\mathbf 1_{I_j}$, $k_{q+1}=+1$ and the labels $\beta_j=-1$, $\beta_{q+1}=q-\varepsilon$. With the single blind cell and equal weights $\lambda_\ell=1/(q+1)$ the pooled kernel is $\frac1{q+1}\bigl(1-\sum_j\mathbf 1_{I_j}\bigr)=0$ because the $I_j$ tile $[0,T]$; hence the kernel term of (6) vanishes and
\[
\Gam(\lambda,a)=-\sum_\ell\lambda_\ell\beta_\ell=-\frac{-q+(q-\varepsilon)}{q+1}=\frac{\varepsilon}{q+1}>0 .
\]
On the aligned mesh with blocks $I_1,\dots,I_q$ the safety-row matrix has rows $-e_j$ ($j\le q$) and $\mathbf 1^\top$ in $\R^q$: $q+1$ rows with exactly one linear relation, so the information--time rank is $r=q$ (and the continuous rank $r_\calI$ is the same, the kernels spanning the $q$-dimensional space of step functions on the $I_j$). Any certificate supported on a proper subset of the modes would, by Proposition~4.3, certify the corresponding proper sub-belief nonviable, contradicting the explicit safe controls; hence every positive certificate has support $q+1=r+1$, and the bound of Theorem~3.2(v) is attained with $m=1$. \qed

\subsection{Remark 4.4 (why Helly does not apply in input dimension)}
Helly's theorem bounds the size of an infeasible subsystem by one plus the dimension of the space in which the constraints are half-spaces. The finite program's variables are the block moments $v\in\R^{mQ}$, and the safety rows are half-spaces in that space, not in $\R^m$: in the construction above, $m=1$ but the rows live in $\R^{q}$ (one block per interval), and the obstruction genuinely needs all $q+1$ of them. The statement ``$m+1$ compatible states suffice'' would be true only if the control were a single vector, i.e.\ if the information structure had one block; it is false as soon as the control is a function of time. \qed
\section{The hypothesis moduli from primitive data}\label{S:moduli}

The constants of Theorem~3.2 are not free assumptions; under the standing bounds they are computed from the data. Let $\|A_\theta\|_\infty\le\bar A$, $\|B_\theta\|_\infty\le\bar B$, $\|E_\theta\|_\infty\le\bar E$, $\|g_\theta\|_\infty\le\bar g$, $\|c_i\|\le\bar c$ on $[0,T]$; $U\subset c_U+R_U\mathbb B$, $D_\theta\subset R_D\mathbb B$, $P_\theta\subset R_P\mathbb B$; $\TV(B_\theta)\le V_B$, $\TV(E_\theta)\le V_E$; and $M_X:=e^{\bar AT}$, so that $\|X_\theta(t,\sigma)\|\le M_X$ (Gr\"onwall on the transition matrix) and $\partial_\sigma X_\theta(t,\sigma)=-X_\theta(t,\sigma)A_\theta(\sigma)$.

\paragraph{Total variation of the kernels (H2), (H4).}
For $k_a(\sigma)=\mathbf 1_{\{\sigma\le t\}}B_\theta(\sigma)^\top X_\theta(t,\sigma)^\top c_i$, the product rule for functions of bounded variation gives on $[0,t]$
\begin{multline*}
\TV(k_a;[0,t])\le\sup_\sigma\|X_\theta(t,\sigma)^\top c_i\|\,\TV(B_\theta)+\sup_\sigma\|B_\theta(\sigma)\|\int_0^t\|X_\theta(t,\sigma)A_\theta(\sigma)^\top c_i\|\dd\sigma\\
\le\bar c\,(M_XV_B+\bar BM_X\bar AT),
\end{multline*}
and the drop to zero at $\sigma=t$ adds at most $\|B_\theta(t)^\top c_i\|\le\bar B\bar c$. Hence $V_U=\bar c(\bar B+M_XV_B+\bar BM_X\bar AT)$, and identically $V_D=\bar c(\bar E+M_XV_E+\bar EM_X\bar AT)$ for the disturbance kernels.

\paragraph{$L^1$ label-continuity modulus (H3).}
For $t\ge t'$, Gr\"onwall on $X_\theta(t,\sigma)-X_\theta(t',\sigma)=\int_{t'}^tA_\theta(r)X_\theta(r,\sigma)\dd r$ gives $\|X_\theta(t,\sigma)-X_\theta(t',\sigma)\|\le M_X\bar A|t-t'|$ uniformly in $\sigma$, so
\[
\|k_t-k_{t'}\|_{L^1}\le\bar B\bar c\,|t-t'|+\int_0^{t'}\bar B\bar cM_X\bar A|t-t'|\dd\sigma\le\bar B\bar c(1+M_X\bar AT)|t-t'| ,
\]
and $|\beta_t-\beta_{t'}|\le L^{\mathrm{tr}}|t-t'|$ by the trajectory bound below. The label map is $L^1$-Lipschitz in $t$ with constant $\bar B\bar c(1+M_X\bar AT)+L^{\mathrm{tr}}$, which is the form the proof of Theorem~3.2(iv) consumes; the sup-norm form is false, because the step kernels jump at $\sigma=t$ for every $t$.

\paragraph{Uniform time-Lipschitz constant (H1).}
For a trajectory $x$ of mode $\theta$, $\|x(t)\|\le M_X\bigl(R_P+T(\bar g+\bar B(R_U+\|c_U\|)+\bar ER_D)\bigr)$ by Gr\"onwall, and
\begin{gather*}
\bigl|\tfrac{\mathrm d}{\mathrm dt}c_i^\top x(t)\bigr|\le\bar c\bigl(\bar A\|x(t)\|+\bar g+\bar B(R_U+\|c_U\|)+\bar ER_D\bigr)\le L^{\mathrm{tr}},\\
L^{\mathrm{tr}}:=\bar c\bar AM_XR_P+\bar c(1+\bar AM_XT)\bigl(\bar g+\bar B(R_U+\|c_U\|)+\bar ER_D\bigr).
\end{gather*}
Each $\calF_{(\theta,i,t)}(\pi)$ is a supremum over $(x_0,d)$ of $L^{\mathrm{tr}}$-Lipschitz functions of $t$, and a pointwise supremum of $L$-Lipschitz functions is $L$-Lipschitz; hence (H1) holds with $L=L^{\mathrm{tr}}$.

\paragraph{The instance.}
For the double integrator with $n=4$, $\bar A=1$, $\bar B=1$, $R_U=1$, $R_P=\tfr{34}{25}$ the closed forms apply but are loose; the direct kernel computation gives the sharp constants used in Section~5.5 of the main text. Position rows have kernels $k(\sigma)=(t-\sigma)n_j\mathbf 1_{\{\sigma\le t\}}$, continuous at $\sigma=t$, with $\TV=t\le T$; velocity rows have kernels $\pm e_i\mathbf 1_{\{\sigma\le t\}}$ with $\TV=1$: $V_U=\max(T,1)=\tfr65$. Since $|v_j(t)|\le1+t\le\tfr{11}5$ along every admissible trajectory and $|\dot v|\le1$, $L=\tfr{11}5$. Scenarios are exact (singleton $P_\theta$, no disturbance): $\delta_\beta=0$. For nonlinear or interval-enclosed data a validated flow-error bound enters at the same two places---it inflates the stored labels $\beta^+_a$ (hence $\delta_\beta$) and the kernel errors $\ebar_a$---and the sandwich (11) is then charged with the inflated budgets; the relaxation direction must be toward a larger feasible controller set, as in Section~8.1 of the main text.

\paragraph{The inner-label sharpening (9) and the bound $Th/4$.}
On a mesh whose label times are block boundaries, a block $J=[s_0,s_1]$ with midpoint $m$ lying before the label time $t$ of a position row has $\bar k-k=(\sigma-m)n_j$ on $J$, and
\[
\int_Jh_U\bigl((\sigma-m)n_j\bigr)\dd\sigma=\int_J|\sigma-m|\,h_U(\pm n_j)\dd\sigma=\tfrac14|J|^2 ,
\]
because $h_U(\pm n_j)=1$. Blocks after $t$ contribute nothing (the kernel vanishes), and velocity rows have piecewise-constant kernels with zero block error. Hence the exact moment error of the label $(j,\text{critical},t)$ is $e_t=\tfrac14\sum_{J\subseteq[0,t]}|J|^2$, which is $th/4\le Th/4$ on the uniform mesh and equals $\tfrac14\sum_J|J|^2\le Th/4$ at the horizon on any aligned mesh with $\max_J|J|=h$. The uniform bound $R_UV_Uh=\tfr65h$ is a factor $4$ larger at the horizon ($Th/4=\tfr3{10}h$). Both are legitimate inputs to program (10); the main text uses the sharp one.

\paragraph{The loose-budget example of Section~5.5.}
Feeding program (10) at $h=\tfr1{10}$ the inflated budget $\ebar_a=e_a+\tfr12$ on every row returns $\rho=\tfr3{100}-\tfr12=-\tfr{47}{100}$ (the dual witness is unchanged; the objective drops by the added budget), and the sandwich (11) gives $J\le-\tfr{47}{100}+\max_a\ebar_a+Lh_t=-\tfr{47}{100}+\tfr{53}{100}+\tfr{11}{100}=\tfr{17}{100}$, still valid. The formula with the uniform special case $R_UV_Uh_u$ in place of $\max_a\ebar_a$ would give $-\tfr{47}{100}+\tfr{12}{100}+\tfr{11}{100}=-\tfr6{25}<\tfr3{50}=J$, which is false: the sandwich must be charged with the budget the program actually used.
\section{Exact arithmetic of the pair policies and of the ladder}\label{S:pairs}

All quantities are for the instance of Section~5 of the main text: $\ddot p=u$, $u\in U$ (the hexagon with vertices $\pm n_1,\pm n_2,\pm n_3$, $h_U(\pm n_j)=1$), $p_j(0)=\tfr{34}{25}n_j$, $v_j(0)=n_j$, facets $n_i^\top p\le2$ and $|v_i|\le\tfr65$, $\tau=\tfr15$, $T=\tfr65$. Along branch $j$ write $\zeta_j(t)=n_j^\top p_j(t)$ and $\nu_j(t)=n_j^\top v_j(t)$ for the critical coordinate and its velocity. Under a constant control $w$ on an interval of length $\ell$ starting from $(\zeta,\nu)$, $\zeta\mapsto\zeta+\nu\ell+\tfrac12(n_j^\top w)\ell^2$ and $\nu\mapsto\nu+(n_j^\top w)\ell$; braking at $u=-n_j$ from velocity $\nu>0$ to rest adds $\nu^2/2$ to $\zeta$ and takes time $\nu$.

\subsection{Proposition 5.2: the two pair policies at $\tau=\tfr15$}

\paragraph{Pair $\{1,2\}$, blind control $w=(-\tfr25,-\tfr45)$.}
$n_1^\top w=n_2^\top w=-\tfr25$ (both in-belief branches brake at the shared authority $\tfr25$). Window exit, both branches: $\nu=1-\tfr25\cdot\tfr15=\tfr{23}{25}$, $\zeta=\tfr{34}{25}+\tfr15-\tfrac12\cdot\tfr25\cdot\tfr1{25}=\tfr{194}{125}$. Braking to rest: peak $\zeta=\tfr{194}{125}+\tfrac12(\tfr{23}{25})^2=\tfr{1940+529}{1250}=\tfr{2469}{1250}=1.9752$, rest at $t=\tfr15+\tfr{23}{25}=\tfr{28}{25}<T$; slack $\tfr{31}{1250}$. Velocity components at the junctions: branch~1: $(1,0)\to(\tfr{23}{25},-\tfr4{25})\to(0,-\tfr4{25})$; branch~2: $(-\tfr35,\tfr45)\to(-\tfr{17}{25},\tfr{16}{25})\to(-\tfr{16}{125},-\tfr{12}{125})$. Every component stays in $[-1,1]$, so the velocity facets $|v_i|\le\tfr65$ hold with margin at least $\tfr15$. Noncritical position facets: for $i\ne j$, $n_i^\top p_j(t)\le(\tfr{34}{25}+t)\,n_i^\top n_j+\tfrac12t^2\le\tfr{18}{25}$ on $[0,T]$ since $n_i^\top n_j\le-\tfr7{25}$ and $h_U(n_i)=1$. The pair $\{1,3\}$ is the reflection $v_2\mapsto-v_2$.

\paragraph{Pair $\{2,3\}$, blind control $w=n_1=(1,0)$.}
$n_2^\top n_1=n_3^\top n_1=-\tfr35$. Window exit: $\nu=1-\tfr35\cdot\tfr15=\tfr{22}{25}$, $\zeta=\tfr{34}{25}+\tfr15-\tfrac12\cdot\tfr35\cdot\tfr1{25}=\tfr{387}{250}$. Braking: peak $\tfr{387}{250}+\tfrac12(\tfr{22}{25})^2=\tfr{1935+484}{1250}=\tfr{2419}{1250}=1.9352$, rest at $t=\tfr{27}{25}$; slack $\tfr{81}{1250}$. Velocity components, branch~2: $(-\tfr35,\tfr45)\to(-\tfr25,\tfr45)\to(-\tfr25,\tfr45)+\tfr{22}{25}(\tfr35,-\tfr45)=(\tfr{16}{125},\tfr{12}{125})$; branch~3 is the mirror image. All components have magnitude at most $\tfr45$.

\paragraph{Falsified variants (recorded so that the policies are not read as generic).}
(i) Holding $u=0$ through the blind window and then braking at $-n_j$ peaks at $\tfr{34}{25}+\tfr15+\tfr12=\tfr{103}{50}>2$ on every branch. (ii) For $\{2,3\}$, inserting a reduced-brake leg $u=-\tfr35n_j$ of length $\tfr15$ after the window before full braking gives $(\zeta,\nu)=(\tfr{387}{250}+\tfr{22}{25}\cdot\tfr15-\tfrac12\cdot\tfr35\cdot\tfr1{25},\ \tfr{22}{25}-\tfr3{25})=(\tfr{214}{125},\tfr{19}{25})$ at $t=\tfr25$ and then peaks at $\tfr{214}{125}+\tfrac12(\tfr{19}{25})^2=\tfr{2501}{1250}>2$ at $t=\tfr{29}{25}$. (iii) The \emph{off-belief} branch~1 under the $\{2,3\}$ blind control $u=n_1$ accelerates to $v_1(\tau)=\tfr65$, exactly the velocity bound, and would then peak at $\tfr{34}{25}+\tfr15+\tfr1{50}+\tfr12(\tfr65)^2=\tfr{23}{10}$: the pair policy is a policy for the pair, not for the triple.

\subsection{Theorem 6.2: junction values along the threshold schedules}

Let $t=t^\ast(B)$, $q=1-t\tau$, and the schedule ``$w$ on $[0,\tau]$; $-n_j$ on $[\tau,\tau+q]$; $0$ on $[\tau+q,T]$''. The critical coordinate obeys
\[
\zeta_j(\tau)=\tfr{34}{25}+\tau-\tfrac t2\tau^2,\qquad\nu_j(\tau)=1-t\tau=q,\qquad\zeta^{\max}=\zeta_j(\tau)+\tfrac12q^2=\tfr{93}{50}+(1-t)\tau+\tfrac{t^2-t}2\tau^2 ,
\]
and $\zeta^{\max}=2$ is the quadratic $6\tau^2-30\tau+7=0$ for $t=\tfr25$ and $6\tau^2-20\tau+7=0$ for $t=\tfr35$ (multiply $2-\zeta^{\max}$ by $50$). The smaller roots are $\tau_{12}=\tfr{15-\sqrt{183}}6$ and $\tau_{23}=\tfr{10-\sqrt{58}}6$, with rational brackets $0.2453<\tau_{12}<0.2454$ and $0.3973<\tau_{23}<0.3974$ obtained by sign changes of the exact quadratics (the verifier evaluates them at the bracket endpoints in rational arithmetic). The label time $s_\ast=\tau+q=1+(1-t)\tau$ is $<\tfr65$ at both roots ($1+\tfr35\cdot0.2454<1.148$; $1+\tfr25\cdot0.3974<1.159$).

\emph{Velocity junction values at the roots.} Each velocity coordinate is affine on each leg, so its extreme values occur at $t=0$, $\tau$, $\tau+q$ and $T$. For $B=\{1,2\}$ with $w=(-\tfr25,-\tfr45)$ at $\tau=\tau_{12}$: branch~1, $(1,0)\to(1-\tfr25\tau,-\tfr45\tau)\to\tau(w+\tfr25n_1)=(0,-\tfr45\tau)$; branch~2, $(-\tfr35,\tfr45)\to(-\tfr35-\tfr25\tau,\tfr45-\tfr45\tau)\to\tau(w+\tfr25n_2)=\tau(-\tfr{16}{25},-\tfr{12}{25})$. For $B=\{2,3\}$ with $w=(1,0)$ at $\tau=\tau_{23}$: branch~2, $(-\tfr35,\tfr45)\to(-\tfr35+\tau,\tfr45)\to\tau(w+\tfr35n_2)=\tau(\tfr{16}{25},\tfr{12}{25})$, and branch~3 its mirror image. With $\tau<\tfr25$ every entry has magnitude at most $1$, so $|v_i|\le1<\tfr65$ throughout, including the final hold leg where the velocity is constant. The noncritical position bound $\tfr{18}{25}$ of the previous subsection is delay-independent.

\emph{Monotonicity of the schedule in $\tau$.} $\zeta^{\max}(\tau)$ is a quadratic in $\tau$ with vertex at $\tau=1/t$, hence increasing on $[0,1/t)\supset[0,\tfr12]$; so the schedule is safe for every $\tau\le\tau_B$, and at $\tau=\tau_B$ exactly it touches the ceiling without crossing it (constraints are closed).

\subsection{Proposition 6.1(ii): the shared authorities}

For a pair $B=\{i,j\}$ the objective $\min_{w\in U}\max(n_i^\top w,n_j^\top w)$ is piecewise linear; its minimum over the hexagon is attained at a vertex of $U$ or on the bisector $\{(n_i-n_j)^\top w=0\}$. For $\{1,2\}$: $n_1-n_2=(\tfr85,-\tfr45)$, bisector $w_2=2w_1$; on it $n_1^\top w=w_1$, and the hexagon's extreme point with $w_1$ minimal is where $|w_2|\le\tfr45$ binds, $w=(-\tfr25,-\tfr45)$, giving $-\tfr25$. At the six vertices $\max(n_1^\top w,n_2^\top w)$ takes the values $1,\tfr35,1,\tfr35,-\tfr7{25},\tfr35$ (at $n_1,-n_1,n_2,-n_2,n_3,-n_3$), none below $-\tfr7{25}$; hence $t^\ast(\{1,2\})=\tfr25$. For $\{2,3\}$: bisector $w_2=0$, $n_2^\top w=-\tfr35w_1$, minimal at the vertex $w=(1,0)=n_1$: $t^\ast=\tfr35$. For the triple, $\sum_j\lambda_jn_j=0$ forces $\max_jn_j^\top w\ge0$ for every $w$, with equality at $w=0$. The pooled normals of the pair duals, $\tfr38n_1+\tfr58n_2=(0,\tfr12)$, $\tfr38n_1+\tfr58n_3=(0,-\tfr12)$ and $\tfr12n_2+\tfr12n_3=(-\tfr35,0)$, have support values $h_U(-\nu)$ equal to $\tfr25$, $\tfr25$, $\tfr35$ respectively ($h_U(e_2)=\tfr45$, $h_U(e_1)=1$)---the same numbers, which is why the two-label duals are tight against the schedules.
\section{The two-stage observer schedule (Proposition 6.4(b))}\label{S:observer}

Information structure: the blind cell $\{1,2,3\}$ on $[0,\tau_1)$, the partition $\{1\}\,|\,\{2,3\}$ on $[\tau_1,\tau)$, full revelation from $\tau=\tfr15$; $\tau_1=\tfr3{25}$. Schedule: common control $w=(\tfr3{20},0)$ on $[0,\tau_1]$; in cell $\{1\}$, $u=-n_1$ from $\tau_1$ to rest; in cell $\{2,3\}$, the shared control $u=n_1$ on $[\tau_1,\tau]$ (the optimal shared brake of the pair, authority $\tfr35$), then $u=-n_j$ to rest, then $u=0$.

\paragraph{Branch 1.}
$n_1^\top w=\tfr3{20}>0$: the common control \emph{accelerates} branch~1 slightly, buying braking authority for the other two. At $\tau_1$: $\nu_1=1+\tfr3{20}\cdot\tfr3{25}=\tfr{509}{500}$, $\zeta_1=\tfr{34}{25}+\tfr3{25}+\tfrac12\cdot\tfr3{20}\cdot\tfr9{625}=\tfr{37027}{25000}$. Braking to rest from $\tau_1$ takes $\tfr{509}{500}<T-\tau_1=\tfr{27}{25}$ and adds $\tfrac12(\tfr{509}{500})^2=\tfr{259081}{500000}$: peak $\zeta_1^{\max}=\tfr{740540+259081}{500000}=\tfr{999621}{500000}=1.999242<2$. Velocity components: $(1,0)\to(\tfr{509}{500},0)\to(0,0)$; the largest magnitude $\tfr{509}{500}<\tfr65$.

\paragraph{Branches 2 and 3.}
$n_j^\top w=-\tfr9{100}$. At $\tau_1$: $\nu_j=1-\tfr9{100}\cdot\tfr3{25}=\tfr{2473}{2500}$, $\zeta_j=\tfr{34}{25}+\tfr3{25}-\tfrac12\cdot\tfr9{100}\cdot\tfr9{625}=\tfr{184919}{125000}$. Shared leg $u=n_1$ of length $\tfr2{25}$ with $n_j^\top n_1=-\tfr35$: $\nu_j\mapsto\tfr{2473}{2500}-\tfr6{125}=\tfr{2353}{2500}$, $\zeta_j\mapsto\tfr{184919}{125000}+\tfr{2473}{2500}\cdot\tfr2{25}-\tfrac12\cdot\tfr35\cdot\tfr4{625}=\tfr{194571}{125000}$. Braking to rest from $\tau$ takes $\tfr{2353}{2500}<1$ and adds $\tfrac12(\tfr{2353}{2500})^2=\tfr{5536609}{12500000}$: peak $\tfr{19457100+5536609}{12500000}=\tfr{24993709}{12500000}\approx1.999497<2$. Velocity components for branch~2: $(-\tfr35,\tfr45)\to(-\tfr35-\tfr9{500},\tfr45)\to(-\tfr35-\tfr9{500}+\tfr2{25},\tfr45)\to$ rest; all magnitudes at most $\tfr45$ (branch~3 is the mirror image). Noncritical position facets are below $\tfr{18}{25}$ by the bound of Section~\ref{S:pairs}.

\paragraph{Why these numbers.}
The common control is chosen on the bisector of branches~2 and~3 ($w_2=0$) so that both brake equally; the single free parameter $w_1$ trades branch~1's acceleration against the others' braking. At $w_1=\tfr3{20}$ the three peaks are $1.99924$ and $1.99950$: nearly balanced, which is what a schedule at the edge of viability should look like. The schedule is one feasible point, not an optimum, and the open gap $[\tfr3{25},\tfr{6-\sqrt{29}}5]$ is a statement about this schedule and the two-label dual jointly; closing it would require either a better schedule (pushing the lower end up) or a dual with more labels (pushing the upper end down). The finite program at $h=\tfr1{100}$ does not improve the upper end: its sweep (Table~\ref{tab:sweep}) first turns positive at $\tau_1=\tfr{13}{100}>\tau_1^{\mathrm{dual}}$, because of the moment error $Th/4=\tfr3{1000}$.

\begin{table}[h]
\centering\footnotesize
\caption{The archived sweep of Figure~4: finite program (10) at $h=\tfr1{100}$ for the two-stage structure, HiGHS optimum $\rho(\tau_1)$ and the eroded triple certificate $\Gamma(\tau_1)$. Every $\rho$ is a lower bound on $J$ for that structure; $\rho\ge\Gamma-Th/4$ throughout, with equality at $\tau_1=\tfr15$.}
\label{tab:sweep}
\begin{tabular}{@{}rrr@{\qquad}rrr@{}}
\toprule
$\tau_1$ & $\rho(\tau_1)$ & $\Gamma(\tau_1)$ & $\tau_1$ & $\rho(\tau_1)$ & $\Gamma(\tau_1)$\\
\midrule
$0.00$ & $-0.067500$ & $-0.105000$ & $0.11$ & $-0.011088$ & $-0.010537$\\
$0.01$ & $-0.067500$ & $-0.096037$ & $0.12$ & $-0.003450$ & $-0.002400$\\
$0.02$ & $-0.067500$ & $-0.087150$ & $0.13$ & $+0.004163$ & $+0.005662$\\
$0.03$ & $-0.067500$ & $-0.078338$ & $0.14$ & $+0.011775$ & $+0.013650$\\
$0.04$ & $-0.065095$ & $-0.069600$ & $0.15$ & $+0.019363$ & $+0.021562$\\
$0.05$ & $-0.057325$ & $-0.060937$ & $0.16$ & $+0.026925$ & $+0.029400$\\
$0.06$ & $-0.049575$ & $-0.052350$ & $0.17$ & $+0.034463$ & $+0.037163$\\
$0.07$ & $-0.041837$ & $-0.043838$ & $0.18$ & $+0.042000$ & $+0.044850$\\
$0.08$ & $-0.034125$ & $-0.035400$ & $0.19$ & $+0.049513$ & $+0.052463$\\
$0.09$ & $-0.026434$ & $-0.027037$ & $0.20$ & $+0.057000$ & $+0.060000$\\
$0.10$ & $-0.018750$ & $-0.018750$ & & &\\
\bottomrule
\end{tabular}
\end{table}

For $\tau_1\le\tfr3{100}$ the program's value is constant at $-0.0675$: there the binding sub-belief is the pair $\{2,3\}$ (branch~1 is revealed almost immediately and is harmless on its own), and $-0.0675$ is the finite value of that pair's program, consistent with the pair's exact slack $-\tfr{81}{1250}=-0.0648$ minus a moment error of at most $\tfr3{1000}$. This is the sense in which the optimised program is more informative than the fixed-weight dual: it re-optimises the weights and labels as the information structure changes, and its value switches from the triple's certificate to the pair's.
\section{Proofs for the extensions (Section 7)}\label{S:extensions}

\subsection{The erosion identity (16)}
$\phi_U(g)=\min_{v\in U}g^\top v$ is concave and positively homogeneous, hence superadditive: $\phi_U(g_1+g_2)\ge\phi_U(g_1)+\phi_U(g_2)$, because a common minimiser for the sum is admissible for each summand separately. Evaluating one fixed witness $(\lambda,a)$ on the coarse structure and on its refinement changes only the pooling on $I$: the kernel term $\int_I\phi_U(\sum_jg_{C_j})$ becomes $\int_I\sum_j\phi_U(g_{C_j})$, and the label term is unchanged. The difference is the displayed integral, which is nonnegative pointwise. For $U=[-\bar u,\bar u]^m$, $\phi_U(g)=-\bar u\|g\|_1$, giving the $\ell^1$ form. The software regression in the verifier (superadditivity on $2025$ deterministic rational pairs for the hexagon) is a check on the implementation of $\phi_U$, not evidence for the identity. \qed

\subsection{The oracle-recourse certificate (17)}
The oracle relaxation replaces the true post-$\tau$ policy class by the larger class that may assign each stored scenario its own measurable control. The expression (17) is the atomic margin (6) evaluated in that enlarged class: before $\tau$ the pooled kernel term is $\phi_U(\sum_\ell\lambda_\ell k_\ell)$ (one common control for all scenarios, which is forced on every real policy because the scenarios generate identical histories), after $\tau$ each scenario $\ell$ contributes $\phi_U(\lambda_\ell k_\ell)$ with its own label $\beta_\ell$. A positive margin certifies, by Lemma~3.1 applied to the enlarged class, that every policy of the enlarged class violates a stored row; every real output-feedback policy belongs to the enlarged class, so none is safe. Completeness fails because the enlarged class may contain safe policies that no real policy can imitate (it knows the scenario after $\tau$ even when histories remain identical). \qed

\subsection{Proposition 7.1 (belief cells)}
Let each cell $C$ carry a verified outer enclosure: for every label $a=(C,i,t)$ a rational $\beta_a^+\ge\sup\{\beta_{(\theta,x_0,d),i,t}:(\theta,x_0)\in C\}$ with $\beta^+_a-\beta_a\le\delta_\beta$. The row functional with the enclosed label dominates the true one for every policy, so the finite program built with the enclosed labels has a feasible set containing the moment image of every safe cell-adapted policy; by the argument of Theorem~3.2(ii) its dual certificates are sound for the cell-filtration class. The lower bound (i) holds with $\beta^+$ in place of $\beta$ and the same $\ebar$. An upper bound needs more: the held policy constructed from the program's optimiser must be implementable with the information the cells actually provide, and its continuous rows must be evaluated with a budget that includes the full label discrepancy---not only the radius of a state cell---because the label is a worst case over the cell of a functional that depends on the mode as well as the state. Hence the sandwich requires the separately verified (H6) budget, as stated.

\emph{Why the label term moves separately under a split.} Consider $\dot x=u$, $u\in[-1,1]$, $T=1$, one row $x(T)\le1$, no observations, and a belief cell consisting of the initial states $x_0\in[-2,\tfr12]$, which is split into $[-2,-1]$ and $[-1,\tfr12]$. The kernel of the row is $k\equiv1$ for every scenario, so for any weights the pooled kernel is additive, $\phi_U(\lambda_1k+\lambda_2k)=-(\lambda_1+\lambda_2)=\phi_U(\lambda_1k)+\phi_U(\lambda_2k)$: the kernel term of (16) is exactly zero. The labels are the cell suprema of $x_0-1$: coarse $-\tfr12$; fine $-2$ and $-\tfr12$. With equal weights the label term $\sum\lambda\beta$ moves from $-\tfr12$ to $-\tfr54$, i.e.\ the margin \emph{rises} by $\tfr34$ under the split although the kernel erosion is zero. The erosion identity accounts for the control-kernel grouping and nothing else; when a refinement changes the labels, the two effects must be tracked separately, and they can have opposite signs.

\emph{Budgets on the instance.} For a proposed noisy state cell around the instance's trajectories, the slacks $\tfr{31}{1250}$ and $\tfr{81}{1250}$ of the pair policies bound the admissible held-policy row discrepancy, and the dual margin $\tfr3{50}$ bounds the total label error chargeable to the triple certificate, $\sum_j\lambda_j(\beta^+_j-\beta_j)<\tfr3{50}$; a position radius alone bounds neither without a verified enclosure of the mode-dependent terms. \qed
\section{The dilation extrapolation and why it is not a threshold law (Section 6.3)}\label{S:dilation}

Replace $U$ by $aU$, $a\ge1$, keep all other data of the instance, and evaluate the triple certificate's \emph{algebraic form} with $h_{aU}=ah_U$: the blind-window term stays zero (the pooled kernel vanishes), the post-window term scales by $a$, and the label term is unchanged, which gives the formal expression
\[
\Gamma_a(\tau)=\tau-\tfr7{50}-\tfrac{a-1}2 ,
\]
with zeros at $\tau=\tfr7{50},\tfr{39}{100},\tfr{16}{25}$ for $a=1,\tfr32,2$. The expression is a correct value of the triple certificate with its original weights and labels evaluated on the dilated program, hence a lower bound on the dilated system's value, and its zeros are therefore \emph{not} thresholds: positivity certifies nonviability, but non-positivity certifies nothing.

The dilated system is in fact nonviable well before the formal zero. At $a=2$ and $\tau=\tfr{16}{25}$ consider the $\lambda$-weighted critical position $\bar\zeta(t)=\sum_j\lambda_j\zeta_j(t)$ along any policy. On the blind window the common control contributes $\sum_j\lambda_jn_j^\top w=0$, so $\bar\zeta(t)=\tfr{34}{25}+t$ regardless of the control, reaching $2$ at $t=\tfr{16}{25}=\tau$ with $\dot{\bar\zeta}=1$. After $\tau$, even if every branch is granted its own control in $2U$ (deceleration at most $2$ along each $n_j$), $\bar\zeta(\tau+\sigma)\ge2+\sigma-\sigma^2$, which equals $\tfr{209}{100}$ at $\sigma=\tfr1{10}$, i.e.\ at $t=\tfr{37}{50}<T$; since a weighted average exceeds $2$, some branch violates its facet. The dilated system at $\tau=\tfr{16}{25}$ is nonviable although $\Gamma_2(\tfr{16}{25})=0$, and the true threshold for $a=2$ lies well below $\tfr{16}{25}$: the same weighted-average argument, with the post-window rise $\tfrac12\cdot1^2/2=\tfr14$ at deceleration $2$, shows nonviability whenever $\tfr{34}{25}+\tau+\tfr14>2$, i.e.\ for every $\tau>\tfr{39}{100}$. The lesson is the one stated in the main text: a certificate's algebraic sensitivity to a design parameter is a property of that certificate, not of the system; redesign claims require a new certificate or a new schedule.
\section{The verification record}\label{S:verification}

All numbers in the main text are re-derived in exact rational arithmetic (Python \texttt{fractions}); floating point enters only inside the solver call, whose output is compared against the exact value. Every script is deterministic, uses the standard library plus NumPy/SciPy/Matplotlib where stated, and is archived with the manuscript at \url{https://github.com/MIKEAA2020/general-sustainability}.

\subsection{Scripts}

\begin{description}[leftmargin=1.5em,itemsep=2pt]
\item[\texttt{p03\_core.py}] The exact module shared by all v23 scripts: instance data; $h_U$, $\phi_U$, membership in $U$; exact leg-wise trajectory arithmetic (\texttt{critical\_peak}: the critical coordinate is quadratic on each constant-control leg, so its maximum is at a leg endpoint or at the interior stationary point, both rational; \texttt{facet\_values}: an exact-arithmetic grid cross-check of all seven facets at step $\tfr1{400}$, used only to corroborate the closed forms of Section~\ref{S:pairs}); the builder of program (10) for an arbitrary aligned mesh and an arbitrary refining information structure (\texttt{build\_lp}); the closed-form primal and dual witnesses (\texttt{witness\_primal}, \texttt{witness\_dual}) and the mesh law; a HiGHS interface (\texttt{solve\_highs}, SciPy \texttt{linprog}, dual simplex); the erosion integrals, shared authorities, pair duals, exact root brackets by rational bisection on the exact quadratics, and exact Gaussian-elimination rank.
\item[\texttt{paper03\_computational\_certification\_v23\_verification.py}] 63 checks in eleven groups: (A) instance geometry, 7 checks; (B) the triple's exact threshold, 6; (C) pair viability and the falsified variants, 8; (D) shared authorities and the ladder, 8; (E) the finite program---dimension identities, witness values at the five campaign meshes, HiGHS agreement, the Farkas contradiction at $h=\tfr1{10}$, the telescoping dual objective, the mesh law on seven further aligned meshes, the smallest certifying aligned program, the coarsest mesh---10; (F) the sandwich on the instance, 6; (G) rank, depth and the scalar-input construction for $q\in\{1,2,3,5,8\}$, 3; (H) the erosion table, 5; (I) the two-stage observer: eroded certificate, root bracket, exact schedule peaks, facet checks, monotonicity, and the archived sweep, 7; (J) the dilation scoping, 2; (K) the residual-bounded test and the saturation of the $r+1$ bound, 2. Output: \path{paper03_computational_certification_v23_verification_output.txt} (``verification: 63/63 checks pass''). The HiGHS group is skipped with a notice if SciPy is absent; no exact check depends on it.
\item[\texttt{make\_figs.py}] Produces Figures~1--4 by recompute-then-assert: every plotted peak, threshold, root bracket and crossing is recomputed from \texttt{p03\_core} and asserted against the printed values before drawing; the observer sweep of Figure~4 (21 HiGHS solves at $h=\tfr1{100}$, about 35~s) is written to \texttt{observer\_sweep.tsv}, which the verifier then reads back (group~I).
\item[\texttt{paper2c\_lp\_campaign.py}] (inherited, 9/9) The original solver campaign at the five uniform meshes with the exact witness pair; its label ``12193/6250'' for the pair $\{2,3\}$ peak is a typographical slip for $\tfr{2419}{1250}$, which the v23 verifier states correctly.
\item[\texttt{paper2\_nonstandard\_verification.py}] (inherited, 51/51) The predecessor exact record: instance values, the $m+1$ counterexample, superadditivity regression, error-bound checks. The v22 description of the pair $\{2,3\}$ policy as a ``two-phase brake $-\tfr25n_j$ then $-n_j$'' was incorrect (such a leg would peak above $2$, cf.\ Section~\ref{S:pairs}); the policies of Proposition~5.2 are the ones verified here.
\item[\texttt{certificate\_exchange\_v1.py}] The independent certificate-file parser, validating archived witnesses of this paper and of the measure-dual companion in either direction.
\item[\texttt{viacert/}] The finite reference library (standard library only): \texttt{python3 -m viacert selftest} runs the twelve audited identities (pair family 36; corner singleton nonviable; four-cell counts $(24,26,25,28)$; meet of middles equals the institutional kernel setwise; partition census $15/7/2$; maximal common-action overlap at $(2,2)$; benchmark $Y^\ast=\tfr{27}5$ with cap sum $2$; CE drift bounds; the non-strict delay identity on 48 cells, its strict failure at the boundary cell $(2,1)$, $\sigma^\ast(2)=1$; step-1 witness extraction), all passing. The package implements the finite calculus and monitoring criteria; it does not reimplement the continuous programs.
\end{description}

Timings in Table~1 are single runs on the verification machine (Linux x86\_64, Intel Xeon @ 2.60\,GHz, Python 3.13, SciPy 1.17, HiGHS dual simplex with default settings).

\subsection{The mechanised layer}

The checks above are computations; they do not prove. Separately, the Farkas soundness direction on which the exact witnesses rest is proved in Lean~4 (60 preserved modules under \texttt{Formalizations/}; the committed project pin \texttt{v4.34.1} rebuilt in GitHub Actions with the default Lake target at 60/60 jobs, exit~0: \url{https://github.com/MIKEAA2020/general-sustainability/actions/runs/36946730855}). A comment-aware scan of all 62 source files finds no \texttt{sorry}, \texttt{admit}, explicit \texttt{axiom} or \texttt{constant} declaration; twenty-one named footprints contain no \texttt{sorryAx}. The aggregate-import environment exposes three source-level \texttt{native\_decide} uses in support/value examples, yielding nine generated axiom names under this pin: executable computational certificates that extend kernel-only trust to the Lean compiler and runtime, declared as such and gated (declared-exceptions gate passed in run 36990482599; the separate strict zero-exception check exits~2 by design, and the two verdicts must not be conflated). Three preserved historical modules outside the aggregate import closure are not certified by the default build. The converse Farkas direction (its classical argument rests on a separation theorem the layer does not develop), the value identity, the bridge theorem and the belief-cell proposition call for analysis an ordered-field interface does not supply and are cited rather than formalised. That is a statement about the formalisation, not about the results.

\subsection{Figure provenance}

Figure~1: the five trajectories (three audited policies, two falsified variants) are the exact leg-wise solutions of Section~\ref{S:pairs}, plotted from their closed forms; the five peaks are asserted equal to $\tfr{93}{50},\tfr{2469}{1250},\tfr{2419}{1250},\tfr{103}{50},\tfr{2501}{1250}$ before drawing. Figure~2: finite values from the mesh law on the six uniform aligned meshes (asserted against the witness pair), upper bounds $\rho+Th/4+Lh/2$ (asserted to equal $\tfr7{25},\tfr{17}{100},\tfr{71}{1000}$ at $h=\tfr15,\tfr1{10},\tfr1{100}$). Figure~3: $\Gamma$ and the two $\Gamma_B$ from their closed forms, roots asserted inside the rational brackets. Figure~4: the eroded certificate from its closed form, the sweep from \texttt{observer\_sweep.tsv}, and the bracket endpoints $\tfr3{25}$ and $\tfr{6-\sqrt{29}}5$.
\section{Sparsification on the instance}\label{S:sparsification}

At $h=\tfr1{10}$ program (10) has $465$ inequalities, of which $Ms(N_t+1)=3\cdot7\cdot13=273$ are safety rows and $p_UQ=6\cdot32=192$ are input-facet rows; the safety-row matrix has full column rank $r=mQ=64$, so Theorem~3.2(v) allows a positive certificate with up to $65$ safety labels. The dual witness of Proposition~5.3 uses three: the horizon labels $(j,\text{critical facet},T)$, $j=1,2,3$, with weights $\lambda=(\tfr38,\tfr5{16},\tfr5{16})$, plus input-facet multipliers $\mu_{j,k}=\lambda_ja_k\eta_j$ on the post-observation blocks (these are not safety labels and do not count toward the support). The deletion-based procedure of Section~8.1 run on this certificate with target margin $\gamma_0=\tfr3{100}$ attempts each of the three supported labels in turn and keeps none of the deletions: removing the label of branch $j$ leaves the labels of the complementary pair, and the pair's finite program is feasible (its optimum is at most the pair's exact slack, $-\tfr{31}{1250}$ or $-\tfr{81}{1250}$, below zero), so no certificate---let alone one with the target margin---survives. The support $\{1,2,3\}$ is therefore irredundant, and by Proposition~4.3 it is minimal, since the obstruction depth of the triple at $\tau=\tfr15$ is $3$. The gap between $3$ and the admissible $65$ is the statement of Corollary~5.5 in computational form: the rank bound is a worst case over all row systems of the given dimensions, and the depth is the instance's own conflict structure.

The same conclusion holds on every aligned mesh with $h<\tfr15$---for instance on the $2^7$-fold refinement of $\{0,\tau,T\}$ with $8469$ rows and up to $mQ+1=1025$ admissible labels---because the witness pair is mesh-generic and the pairs' programs remain feasible: the certificate's support does not grow with the mesh, only the program does. This is the practical content of the size result: on a nonviable instance one may expect the \emph{certificate} to stay small as the \emph{program} is refined, and the residual-bounded test (18) checks it in time linear in the support.

\end{document}
